\documentclass[11pt]{amsart}
\usepackage{amsmath,amssymb,amsthm}
\usepackage[margin=1.05in]{geometry}
\raggedbottom

\numberwithin{equation}{section}
\newtheorem{theorem}{Theorem}[section]
\newtheorem{proposition}[theorem]{Proposition}
\newtheorem{lemma}[theorem]{Lemma}
\newtheorem{corollary}[theorem]{Corollary}
\theoremstyle{definition}
\newtheorem{definition}[theorem]{Definition}
\theoremstyle{remark}
\newtheorem{remark}[theorem]{Remark}
\newtheorem{example}[theorem]{Example}

\newcommand{\Z}{\mathbb{Z}}
\newcommand{\Q}{\mathbb{Q}}
\newcommand{\R}{\mathbb{R}}
\newcommand{\N}{\mathbb{N}}
\newcommand{\C}{\mathbb{C}}
\newcommand{\F}{\mathbb{F}}
\newcommand{\T}{\mathbb{T}}
\newcommand{\hZ}{\widehat{\Z}}
\newcommand{\cA}{\mathcal{A}}
\newcommand{\cO}{\mathcal{O}}
\newcommand{\cP}{\mathcal{P}}
\newcommand{\Ns}{\mathrm{N}}
\newcommand{\Pic}{\operatorname{Pic}}
\newcommand{\Jac}{\operatorname{Jac}}
\newcommand{\End}{\operatorname{End}}
\newcommand{\Div}{\operatorname{Div}}
\newcommand{\Gal}{\operatorname{Gal}}
\newcommand{\Idl}{\mathbb{I}}
\newcommand{\tors}{\mathrm{tors}}
\newcommand{\id}{\mathrm{id}}

\PassOptionsToPackage{hyphens}{url}
\usepackage[hidelinks,bookmarks=false]{hyperref}

\title[The limits of abelian reconstruction]{The limits of abelian reconstruction:\\ twin curves exist, genus one suffices,\\ and the obstruction is a torsor coordinate}
\author{Ruqing Chen}
\address{GUT Geoservice Inc., Montreal, Canada}
\email{ruqing@hotmail.com}
\thanks{Sources, code and data for the entire series: \url{https://github.com/Ruqing1963/localized-bost-connes}; this paper is in \texttt{papers/XX-twin-curves/}.}
\date{}

\begin{document}

\begin{abstract}
Paper XIX of the series asked whether there are non-isomorphic curves over $\F_q$ with the
same zeta function \emph{and} isomorphic Jacobian groups, which would locate the resolution
limit of the abelian Bost--Connes systems. The answer is yes, the construction is far simpler
than the genus-two machinery proposed, and the reason is complex multiplication.

For an ordinary elliptic curve $E/\F_q$ with $\End(E)=\cO$, Deuring's theory gives
$E(\F_{q^n})\cong\cO/(\phi^n-1)$ as an $\cO$-module for every $n$, an isomorphism depending
only on $\cO$; and $\Pic(\cO)$ acts simply transitively on the set of curves in the isogeny
class with that endomorphism ring, of which there are exactly $h(\cO)$. Hence whenever
$h(\cO)>1$ one obtains $h(\cO)$ curves with distinct $j$-invariants sharing \emph{every}
Galois-module invariant at every level. Genus one already suffices; Howe's and Lauter's
genus-two constructions are not needed. Enumerating all curves over $\F_p$ for $11\le p\le29$ we
find $69$ isogeny classes containing several $j$-invariants with a common group structure,
and the counts match the class numbers: $h(-88)=h(-100)=h(-48)=h(-32)=h(-60)=2$ against two
$j$-invariants in each case.

This is stronger than the question asked. Twin curves agree not only in $Z_C$ and in
$\Jac(\F_q)$, but in the whole tower $\{E(\F_{q^n})\}_n$ with its trace maps --- which matters
because Tier 2 of Paper XIX recovers more than the abstract group: the compactified Picard
group with its translation structure, through the deep stratum, and with it the map
$P\mapsto[P]\in\Pic(C)$ on all closed points. We prove that this map does not separate twins
either --- under the module identification the class of a degree-$n$ point is multiplication
by $(\phi^n-1)/(\phi-1)$, again $\cO$-determined --- and confirm it numerically: for two
supersingular classes at $p=11,17$ and for each group structure at $p=13$, the degree-two
point-to-Picard multiset profiles coincide across distinct $j$-invariants.

We also correct the proposed explanation of the barrier. The Weil pairing is not out of
abelian reach for the reason given: over $\F_q$ the Galois group acting on $\Jac[n]$ is
$\Gal(\bar\F_q/\F_q)=\hZ$, which is abelian, so ``the pairing cannot be encoded in an
abelianized Galois group'' is not the obstruction. The obstruction is a torsor coordinate:
$\Pic(\cO)$ acts simply transitively on the twins and acts \emph{trivially} on every module
invariant, so what is missing is the position in a $\Pic(\cO)$-torsor, and no invariant of
the modules can see it.
\end{abstract}

\maketitle

\section{Twin curves, and why genus one suffices}

The framework is that of \cite{Main} transported to function fields as in Papers XVII--XIX
of the series; Papers II, IX, X, XIV and XVI--XIX, referred to by numeral, are unpublished
companions, and every fact used from them is restated or verified in place.

\begin{theorem}[Existence of twins]\label{thm:twins}
Let $E/\F_q$ be an ordinary elliptic curve with Frobenius $\phi$ and $\End(E)=\cO$, an order
in $\Q(\phi)$ with $\Z[\phi]\subseteq\cO\subseteq\cO_{\max}$. Then:
\begin{enumerate}
\item $E(\F_{q^n})\cong\cO/(\phi^n-1)$ as an $\cO$-module, for every $n\ge1$, and this
depends only on $\cO$ and $\phi$, not on $E$;
\item $\Pic(\cO)$ acts simply transitively on the set of $\F_q$-isomorphism classes of
curves in the isogeny class with endomorphism ring $\cO$ --- equivalently, since the
quadratic twist has trace $-t\ne t$ for ordinary $E$, on their $j$-invariants; there are
exactly $h(\cO)$ of them.
\end{enumerate}
Consequently, whenever $h(\cO)>1$ there exist $C_1\not\cong C_2$ over $\F_q$ with
$Z_{C_1}=Z_{C_2}$, with $\Jac(C_1)(\F_q)\cong\Jac(C_2)(\F_q)$, and indeed with
$C_1(\F_{q^n})$ and $C_2(\F_{q^n})$ isomorphic as $\cO$-modules for every $n$.
\end{theorem}

\begin{proof}
(1) is Deuring's theory \cite{Deuring} in the module-theoretic form proved by Lenstra
\cite{Lenstra}: for ordinary $E$ with $\End(E)=\cO$ one has
$E(\F_{q^n})=\ker(\phi^n-1)\cong\cO/(\phi^n-1)$ as $\cO$-modules. (2) is the theory of
complex multiplication in characteristic $p$, in Waterhouse's kernel-ideal form
\cite{Waterhouse}: the set of curves with a given ordinary endomorphism ring is a principal
homogeneous space under the class group. Distinct classes give distinct $j$-invariants, these
being the roots of the Hilbert class polynomial of $\cO$ reduced mod $p$.
\end{proof}

\begin{remark}[The proposal over-engineers]\label{rem:overengineer}
The construction requested --- genus-two hyperelliptic curves with equal Igusa invariants
data, following Howe \cite{Howe} and Lauter \cite{Lauter} --- is unnecessary. Genus one produces twins as soon as
$h(\cO)>1$, which is the generic situation, and produces them with a stronger property than
asked for: agreement of the entire Galois-module tower, not merely of $Z_C$ and
$\Jac(\F_q)$.
\end{remark}

\begin{example}[Enumeration]\label{ex:enum}
Enumerating all $y^2=x^3+ax+b$ over $\F_p$ for $11\le p\le29$ and grouping by (trace, group
structure) yields $69$ classes containing more than one $j$-invariant. Samples, with the
class number of the relevant order:
\[
\begin{array}{lllll}
p & t & \Jac(\F_p) & j\text{-invariants} & h(\cO)\\\hline
23 & -2 & \Z/26 & 1,\ 17 & h(-88)=2\\
13 & 2 & \Z/12 & 4,\ 6 & h(-48)=2\\
17 & -6 & \Z/24 & 1,\ 6 & h(-32)=2\\
19 & -4 & \Z/24 & 9,\ 17 & h(-60)=2\\
29 & -4 & \Z/34 & 3,\ 7,\ 17 & h(-100)=2\ \text{(two orders)}\\
23 & -6 & \Z/30 & 14,\ 18,\ 20,\ 22 & h(-56)=4
\end{array}
\]
The count of $j$-invariants at a fixed group structure equals $h(\cO)$ for the corresponding
endomorphism ring, as Theorem~\ref{thm:twins}(2) requires.
\end{example}

\section{Indistinguishability at both tiers}

\begin{theorem}[Twins are invisible to the abelian system]\label{thm:invisible}
Let $C_1,C_2$ be twins in the sense of Theorem~\ref{thm:twins}, with a common endomorphism
ring $\cO$. Then their Bost--Connes systems agree at both tiers of Theorem 4.1 of Paper XIX:
\begin{enumerate}
\item \emph{Tier 1.} The partition functions coincide, $Z_{C_1}=Z_{C_2}$.
\item \emph{Tier 2.} $\Jac(C_1)(\F_q)\cong\Jac(C_2)(\F_q)$; more, the whole tower
$\{C_i(\F_{q^n})\}_n$ with its norm maps is isomorphic, so every invariant of the
Galois-module data agrees; and more still, the full Tier-2 datum --- the compactified Picard
group together with the degree-labelled multiset of point classes $\{[P]:\deg P=n\}$ ---
corresponds under the module identification, for every $n$.
\end{enumerate}
\end{theorem}

\begin{proof}
Both follow from Theorem~\ref{thm:twins}(1): the tower and its maps are
$\{\cO/(\phi^n-1)\}$ with the natural maps, depending only on $(\cO,\phi)$. For the last
clause of (2): a closed point of degree $n$ is a Frobenius orbit
$\{Q,\phi Q,\dots,\phi^{n-1}Q\}$ of a point $Q\in C(\F_{q^n})$ of exact degree $n$, and,
choosing the rational base point, its class in $\Pic^0=E(\F_q)$ is
$Q\oplus\phi Q\oplus\cdots\oplus\phi^{n-1}Q=\nu_n\cdot Q$ with
$\nu_n=(\phi^n-1)/(\phi-1)\in\cO$, computed in the $\cO$-module $\cO/(\phi^n-1)$, with
$E(\F_q)=\ker(\phi-1)=\nu_n\cdot\cO/(\phi^n-1)$. The degree-$n$ multiset of point classes is
therefore the image multiset of $\nu_n$ on the exact-degree-$n$ elements of
$\cO/(\phi^n-1)$ --- a formula in $(\cO,\phi)$ alone. Since the Tier-2 datum recovered in
Theorem 3.1 of Paper XIX is exactly the compactified Picard group with these degree-labelled
translations, the twins' data are isomorphic as marked groups.
\end{proof}

\begin{remark}[Tier 2 recovers more than the abstract group]\label{rem:tier2more}
It is worth being careful here, because the question as posed assumed Tier 2 sees only
$\Jac(\F_q)$. It sees more. By Theorem 2.1 of Paper XIX, $G_S/\overline{U_S}\cong\widehat{\Pic(C)}$,
and the deep-stratum translation structure recovered in the proof of Theorem 3.1 of Paper
XIX --- the generator $P\in J_S$ acting on the Picard torsor by $\sigma_P$ --- gives the map
\[
P\ \longmapsto\ [P]\in\Pic(C)
\]
on all closed points, together with their degrees. That is
substantially finer than the group $\Jac(\F_q)$, and it is what has to be checked. By
Theorem~\ref{thm:invisible}(2) it is also $\cO$-determined, so twins are not separated by it
either.
\end{remark}

\begin{remark}[Numerical check of the point-to-Picard map]\label{rem:numcheck}
For an elliptic curve with base point $O$, a closed point of degree two is a Galois orbit
$\{Q,\phi Q\}$ in $E(\F_{q^2})$ and its class in $\Pic^0=E(\F_q)$ is $\mathrm{Tr}(Q)=Q\oplus\phi Q$.
The last clause of Theorem~\ref{thm:invisible}(2) predicts equal profiles across twins; we
confirm it by direct computation, for the ordinary twins of Example~\ref{ex:enum} at $p=13$
and also for two supersingular same-zeta-same-group classes (outside the ordinary hypothesis
of Theorem~\ref{thm:twins}, so not covered by the proof):
\[
\begin{array}{llll}
p & t & j & \text{degree-two profile}\\\hline
11 & 0 & 0 & (6^6,5^6)\\
11 & 0 & 1 & (6^6,5^6)\\
17 & 0 & 0 & (9^9,8^9)\\
17 & 0 & 8 & (9^9,8^9)\\
13 & 2 & 4,\ 6 & (8^6,7^6)\\
13 & 2 & 0,\ 11 & (8^9,6^3)
\end{array}
\]
The profiles agree within each twin pair. At $p=13$ the four curves of trace $2$ split into
two profiles, but that split tracks the group structure --- $\Z/12$ for $j=4,6$ and
$\Z/2\times\Z/6$ for $j=0,11$ --- and not the $j$-invariant, so it is Tier 2 information
already accounted for.
\end{remark}

\section{The obstruction is a torsor coordinate, not a pairing}

\begin{remark}[The proposed explanation is not correct]\label{rem:pairing}
The barrier was attributed to the Weil pairing being ``a non-linear skew-symmetric pairing
that cannot be encoded in the abelianized Galois group $\Idl_K/\overline{K^\times}$''. This
reasoning does not work. Over a finite field the Galois group acting on $\Jac(C)[n]$ is
$\Gal(\bar\F_q/\F_q)=\hZ$, which is \emph{abelian}; the Weil pairing is
$\Gal$-equivariant for that action and is not inaccessible on grounds of non-commutativity.
The Tate module with its Frobenius action is precisely the kind of datum an abelian theory
can carry --- indeed $Z_C$ is its characteristic polynomial.
\end{remark}

\begin{theorem}[The correct obstruction]\label{thm:torsor}
Let $\mathcal E(\cO)$ be the set of curves in a fixed ordinary isogeny class with
$\End=\cO$. Then $\Pic(\cO)$ acts simply transitively on $\mathcal E(\cO)$, and acts
\emph{trivially} on every invariant of the Galois-module tower $\{E(\F_{q^n})\}_n$, these being
$\cO/(\phi^n-1)$ independently of the curve. Consequently the missing datum separating twins
is the position of the curve in a $\Pic(\cO)$-torsor, and no invariant constructed from the
modules --- including the Weil pairing, which is likewise $\cO$-determined --- can detect it.
\end{theorem}

\begin{proof}
Simple transitivity is Theorem~\ref{thm:twins}(2). Triviality on module invariants is
Theorem~\ref{thm:twins}(1). A function on $\mathcal E(\cO)$ factoring through the module data
is therefore constant, whereas $j$ is injective on $\mathcal E(\cO)$.
\end{proof}

\begin{remark}[What would be needed]\label{rem:needed}
By Theorem~\ref{thm:torsor} the separation of twins is equivalent to locating a point in a
$\Pic(\cO)$-torsor, and the classical object that does this is the $j$-invariant, i.e.\ the
Hilbert class field of $\cO$ and its Galois action. That is still abelian class field theory
--- but for the \emph{CM field} $\Q(\phi)$, not for the function field $K=\F_q(C)$. The two
abelian theories are different, and the Bost--Connes system of $C$ carries the second, not the
first. This is a sharper diagnosis than ``the barrier is non-abelian'', and it suggests
where a refinement would have to come from: a system built over the CM order rather than over
the curve.
\end{remark}

\section{The resolution boundary}

\begin{theorem}[Master statement]\label{thm:master}
For the localized Bost--Connes system of a curve over $\F_q$:
\begin{itemize}
\item the partition function determines exactly the isogeny class of $\Jac(C)$
(Theorem 4.1 of Paper XVIII);
\item the Cartan pair determines, through the deep stratum, the compactified Picard group
with the degree-labelled point classes $P\mapsto[P]\in\Pic(C)$,
hence $\Jac(C)(\F_q)$ and the rational effective cone, strictly refining the previous line
(Theorem 3.1 of Paper XIX);
\item neither determines the curve: by Theorems~\ref{thm:twins} and \ref{thm:invisible},
whenever $h(\cO)>1$ there are $h(\cO)$ curves agreeing in all of the above and differing in
$j$.
\end{itemize}
The resolution boundary of the abelian theory is therefore the $\cO$-module tower
$\{E(\F_{q^n})\}$, and the residual ambiguity is a $\Pic(\cO)$-torsor of size $h(\cO)$.
\end{theorem}

\section{Retrospective}

\[
\begin{array}{lll}
\text{Question} & \text{Answer} & \text{Reference}\\\hline
\text{do twins exist?} & \textbf{yes, abundantly} & \text{Thm.~1.1, Ex.~1.3}\\
\text{is genus two needed?} & \textbf{no: genus one suffices} & \text{Rem.~1.2}\\
\text{are twins Tier-1 indistinguishable?} & \text{yes} & \text{Thm.~2.1}\\
\text{Tier 2 sees only }\Jac(\F_q)? & \textbf{no: also }P\mapsto[P] & \text{Rem.~2.2}\\
\text{are twins Tier-2 indistinguishable?} & \text{yes, proved and checked} & \text{Thm.~2.1, Rem.~2.3}\\
\text{is the Weil pairing the obstruction?} & \textbf{no; }\Gal(\F_q)\text{ is abelian} & \text{Rem.~3.1}\\
\text{what is the obstruction?} & \textbf{a }\Pic(\cO)\text{-torsor coordinate} & \text{Thm.~3.2}\\
\text{resolution boundary} & \{E(\F_{q^n})\}\ \text{as }\cO\text{-modules} & \text{Thm.~4.1}
\end{array}
\]

The twenty papers have converged on a single principle, and this one states its sharpest
form. The arithmetic that a Bost--Connes system carries is the arithmetic of its equilibrium
states, and that arithmetic is exactly the Galois-module data of the base: the isogeny class
at the thermodynamic level, the Picard group at the level of the Cartan pair. What lies
beyond is not a deeper analytic invariant but a torsor, and a torsor has no invariants at
all --- its points are distinguished only by a coordinate that must be supplied from outside.

That is the honest form of the negative results accumulated here. The $K$-theory did not fail
for lack of computation (Remark 4.5 of Paper II); the von Neumann algebra did not fail for
lack of rigidity technique (Theorem 2.1 of Paper IX); the quantum transport inequalities did not
fail for lack of a better Dirichlet form (Theorem 4.1 of Paper XVI). In each case the invariant
was blind to something that is not an invariant: a measure class, a tail event, and now a
torsor coordinate. The one place where the pattern broke was
Theorem 4.1 of Paper XVII, where commensurability of the log-norms let the factor retain the
temperature --- and that too was arithmetic, not technique.

Three questions we would still pursue. Is the analogue of Theorem~\ref{thm:torsor} true in
higher genus, with $\Pic(\cO)$ replaced by the relevant class group of the CM order, so that
the resolution boundary is uniformly a torsor? Second, Remark~\ref{rem:needed} suggests
building a Bost--Connes system over the CM order $\cO$ rather than over $C$, whose symmetry
group would be $\Pic(\cO)$ acting on the twins; does that system exist and does it separate
them? Third, the twins of Theorem~\ref{thm:twins} are ordinary; the supersingular case has
endomorphisms in a quaternion order, which is where Paper X located the only viable
non-commutative model, and whether the two occurrences of quaternion orders in this programme
are related we do not know.

\begin{thebibliography}{9}
\bibitem{Deuring} M. Deuring, \emph{Die Typen der Multiplikatorenringe elliptischer Funktionenk\"orper}, Abh. Math. Sem. Univ. Hamburg 14 (1941), 197--272.
\bibitem{Howe} E. W. Howe, \emph{Principally polarized ordinary abelian varieties over finite fields}, Trans. Amer. Math. Soc. 347 (1995), 2361--2401.
\bibitem{Lauter} K. Lauter, \emph{The maximum or minimum number of rational points on genus three curves over finite fields}, Compositio Math. 134 (2002), 87--111.
\bibitem{Lenstra} H. W. Lenstra, Jr., \emph{Complex multiplication structure of elliptic curves}, J. Number Theory 56 (1996), 227--241.
\bibitem{Waterhouse} W. C. Waterhouse, \emph{Abelian varieties over finite fields}, Ann. Sci. \'Ec. Norm. Sup. (4) 2 (1969), 521--560.
\bibitem{Main} R. Chen, \emph{KMS state uniqueness and phase transitions in localized Bost--Connes systems}, preprint, 2026, \newline\url{https://doi.org/10.5281/zenodo.22152101}.
\end{thebibliography}

\end{document}
